\subsubsection{Pseudo-monotone Operators}\label{subsec:pseudo-monotone-operators}

In this section we examine a somewhat more general class of nonlinear mappings, called \emph{pseudo-monotone operators}. In applications, it often occurs that the hypotheses imposed in the previous section are unnecessarily strong. In particular, the monotonicity assumption \ref{hyp:monotone-pde-monotonicity} involves both the first-order derivatives and the function itself. As we shall see in this chapter, it is really only necessary to have a monotonicity assumption on the highest-order derivatives: Compactness will take care of the lower-order terms.

\begin{Definition}
\label{def:pseudo-monotone-operator}
Let $X$ be a reflexive Banach space. An operator $T:X\to X^*$ is called \textbf{pseudo-monotone} if $T$ is bounded and if whenever
\begin{equation}\label{eq:pseudo-monotonicity-weak-sequence}
u_j\rightharpoonup\overline u\quad\text{in }X
\end{equation}
and
\begin{equation}\label{eq:pseudo-monotonicity-limsup-hypothesis}
\limsup_{j\to\infty}(T(u_j),u_j-\overline u)\leq0,
\end{equation}
it follows that
\begin{equation}\label{eq:pseudo-monotonicity-liminf-conclusion}
\liminf_{j\to\infty}(T(u_j),u_j-v)\geq(T(\overline u),\overline u-v)\quad\text{for all }v\in X.
\end{equation}

\end{Definition}

The following can be proved using only a slight modification of the proof of the Browder-Minty theorem.

\begin{Theorem}
\label{thm:pseudo-monotone-surjectivity}
Let $X$ be a real reflexive Banach space and suppose $T:X\to X^*$ is continuous, coercive and pseudo-monotone. Then for every $g\in X^*$ there exists a solution $u\in X$ of the equation
\begin{equation}\label{eq:pseudo-monotone-surjectivity-equation}
T(u)=g.
\end{equation}

\end{Theorem}

The proof is left to the reader (Problem~\ref{ex:pseudo-monotone-surjectivity-proof}).

In practice, the following condition is easier to verify than pseudo-monotonicity.

\begin{Definition}
\label{def:calculus-variations-type-operator}
Let $X$ be a reflexive Banach space. An operator $T:X\to X^*$ is said to be \textbf{of the calculus of variations type} if it is bounded, and it has the representation
\begin{equation}\label{eq:calculus-variations-type-diagonal-representation}
T(u)=\widehat T(u,u)
\end{equation}
where the mapping
\begin{equation}\label{eq:calculus-variations-type-two-argument-mapping}
X\times X\ni(u,v)\mapsto\widehat T(u,v)\in X^*
\end{equation}
satisfies the following hypotheses.\begin{enumerate}[label=CV-\arabic*.,ref=CV-\arabic*]\item\label{hyp:cv-diagonal-monotonicity} For each $u\in X$, the mapping $v\mapsto\widehat T(u,v)$ is bounded and continuous from $X$ to $X^*$, and
\begin{equation}\label{eq:calculus-variations-type-diagonal-monotonicity}
(\widehat T(u,u)-\widehat T(u,v),u-v)\geq0\quad\text{for all }v\in X.
\end{equation}
\item\label{hyp:cv-first-variable-continuity} For each $v\in X$, the mapping $u\mapsto\widehat T(u,v)$ is bounded and continuous from $X$ to $X^*$.\item\label{hyp:cv-vanishing-defect-weak-continuity} If
\begin{equation}\label{eq:calculus-variations-type-cv-three-weak-sequence}
u_j\rightharpoonup\overline u\text{ in }X
\end{equation}
and
\begin{equation}\label{eq:calculus-variations-type-vanishing-diagonal-pairing}
(\widehat T(u_j,u_j)-\widehat T(u_j,\overline u),u_j-\overline u)\to0,
\end{equation}
then for every $v\in X$
\begin{equation}\label{eq:calculus-variations-type-weak-image-continuity}
\widehat T(u_j,v)\rightharpoonup\widehat T(\overline u,v)\text{ in }X^*.
\end{equation}
\item\label{hyp:cv-weak-duality-pairing} If
\begin{equation}\label{eq:calculus-variations-type-cv-four-weak-sequence}
u_j\rightharpoonup\overline u\text{ in }X
\end{equation}
and
\begin{equation}\label{eq:calculus-variations-type-cv-four-weak-image}
\widehat T(u_j,v)\rightharpoonup\psi\text{ in }X^*,
\end{equation}
then
\begin{equation}\label{eq:calculus-variations-type-cv-four-pairing-convergence}
(\widehat T(u_j,v),u_j)\to(\psi,\overline u).
\end{equation}
\end{enumerate}
\end{Definition}

As we indicated above, we have the following.

\begin{Theorem}
\label{thm:cv-type-implies-pseudo-monotonicity}
If $T$ is of the calculus of variations type, then $T$ is pseudo-monotone.
\end{Theorem}

\begingroup
\renewcommand{\endProof}{\par\nobreak\hfill$\blacksquare$\par\vspace{0.1cm}}
\begin{Proof}
Let $u_j\rightharpoonup\overline u$ in $X$ and suppose
\begin{equation}\label{eq:cv-pseudo-proof-limsup-assumption}
\limsup_{j\to\infty}(T(u_j),u_j-\overline u)\leq0.
\end{equation}
We wish to show that
\begin{equation}\label{eq:cv-pseudo-proof-liminf-goal}
\liminf_{j\to\infty}(T(u_j),u_j-v)\geq(T(\overline u),\overline u-v)\quad\text{for every }v\in X.
\end{equation}
Since $\widehat T(u_j,\overline u)$ is bounded in $X^*$, we can extract a subsequence $u_j$ such that
\begin{equation}\label{eq:cv-pseudo-proof-bounded-weak-image-subsequence}
\widehat T(u_j,\overline u)\rightharpoonup\psi\quad\text{in }X^*,
\end{equation}
for some $\psi\in X^*$. We now use \ref{hyp:cv-weak-duality-pairing} to get
\begin{equation}\label{eq:cv-pseudo-proof-duality-pair-limit}
\lim_{j\to\infty}(\widehat T(u_j,\overline u),u_j)=(\psi,\overline u).
\end{equation}
Thus, if we define
\begin{equation}\label{eq:cv-pseudo-proof-nonnegative-defect-definition}
x_j:=(\widehat T(u_j,u_j)-\widehat T(u_j,\overline u),u_j-\overline u)\in\mathbb{R},
\end{equation}
we have
\begin{equation}\label{eq:cv-pseudo-proof-defect-limsup-bound}
\begin{aligned}\limsup_{j\to\infty}x_j=\limsup_{j\to\infty}\bigl[&(T(u_j),u_j-\overline u)-(\widehat T(u_j,\overline u),u_j)\\&+(\widehat T(u_j,\overline u),\overline u)\bigr]\leq0.\end{aligned}
\end{equation}
Here we have used \eqref{eq:cv-pseudo-proof-limsup-assumption}, \eqref{eq:cv-pseudo-proof-bounded-weak-image-subsequence} and \eqref{eq:cv-pseudo-proof-duality-pair-limit}. Since \ref{hyp:cv-diagonal-monotonicity} implies that $x_j\geq0$ we get
\begin{equation}\label{eq:cv-pseudo-proof-defect-vanishing}
\lim_{j\to\infty}x_j=0.
\end{equation}
Thus, we can use \ref{hyp:cv-vanishing-defect-weak-continuity} to get
\begin{equation}\label{eq:cv-pseudo-proof-test-image-limit}
\widehat T(u_j,v)\rightharpoonup\widehat T(\overline u,v)\text{ in }X^*\quad\text{for all }v\in X.
\end{equation}
Hence, we can use \ref{hyp:cv-weak-duality-pairing} again to get\[(\widehat T(u_j,v),u_j)\to(\widehat T(\overline u,v),\overline u)\]or
\begin{equation}\label{eq:cv-pseudo-proof-vanishing-test-pairing}
(\widehat T(u_j,v),u_j-\overline u)\to0\quad\text{for all }v\in X.
\end{equation}
We now use this and the fact that $x_j\geq0$ to get
\begin{equation}\label{eq:cv-pseudo-proof-diagonal-comparison}
(T(u_j),u_j-\overline u)\geq(\widehat T(u_j,\overline u),u_j-\overline u)\to0.
\end{equation}
Together with \eqref{eq:cv-pseudo-proof-limsup-assumption} this gives us
\begin{equation}\label{eq:cv-pseudo-proof-diagonal-pairing-limit}
(T(u_j),u_j-\overline u)\to0.
\end{equation}
We now take the inequality\[(T(u_j)-\widehat T(u_j,w),u_j-w)\geq0\quad\text{for all }w\in X\]from \ref{hyp:cv-diagonal-monotonicity}, and plug in
\begin{equation}\label{eq:cv-pseudo-proof-intermediate-test-point}
w=(1-\theta)\overline u+\theta v,
\end{equation}
for $\theta\in(0,1)$. This yields
\begin{equation}\label{eq:cv-pseudo-proof-theta-inequality}
\begin{aligned}\theta(T(u_j),\overline u-v)\geq&-(T(u_j),u_j-\overline u)+(\widehat T(u_j,w),u_j-\overline u)\\&+\theta(\widehat T(u_j,w),\overline u-v).\end{aligned}
\end{equation}
Dividing this by $\theta$ and using \eqref{eq:cv-pseudo-proof-vanishing-test-pairing} and \eqref{eq:cv-pseudo-proof-diagonal-pairing-limit} we get\[\begin{aligned}
\liminf_{j\to\infty}(T(u_j),u_j-v)&=\liminf_{j\to\infty}(T(u_j),u_j-\overline u)+\liminf_{j\to\infty}(T(u_j),\overline u-v)\\
&\geq\liminf_{j\to\infty}(\widehat T(u_j,w),\overline u-v)\\
&=(\widehat T(\overline u,w),\overline u-v)\\
&=(\widehat T(\overline u,(1-\theta)\overline u+\theta v),\overline u-v).
\end{aligned}\]Letting $\theta\searrow0$ we get
\begin{equation}\label{eq:cv-pseudo-proof-liminf-conclusion}
\liminf_{j\to\infty}(T(u_j),u_j-v)\geq(T(\overline u),\overline u-v)\quad\text{for all }v\in X.
\end{equation}
Since this argument holds for any subsequence of the original sequence, the inequality \eqref{eq:cv-pseudo-proof-liminf-conclusion} holds for the entire original sequence. This completes the proof.
\end{Proof}
\endgroup

The following is immediate from the preceding results.

\begin{Corollary}
\label{cor:cv-type-surjectivity}
Let $X$ be a real reflexive Banach space and suppose $T:X\to X^*$ is continuous, coercive and of the calculus of variations type. Then for every $g\in X^*$ there exists a solution $u\in X$ of the equation
\begin{equation}\label{eq:calculus-variations-type-surjectivity-equation}
T(u)=g.
\end{equation}

\end{Corollary}

\subsubsection{Application to PDEs}\label{subsec:cv-type-application-to-pdes}

Let $\Omega\subset\mathbb{R}^n$ be a bounded domain with smooth boundary. We consider quasilinear second-order differential operators having the form
\begin{equation}\label{eq:cv-quasilinear-second-order-operator}
\widetilde A(u)(\bm{x})=-\sum_{i=1}^n\dfrac{\partial}{\partial x_i}a_i(\bm{x},u(\bm{x}),\nabla u(\bm{x}))+a_0(\bm{x},u(\bm{x}),\nabla u(\bm{x})).
\end{equation}
Our goal is to solve the Dirichlet problem for
\begin{equation}\label{eq:cv-quasilinear-dirichlet-equation}
\widetilde A(u)=f
\end{equation}
for appropriate $f$. Formally, we define the bivariate form
\begin{equation}\label{eq:cv-quasilinear-bivariate-form}
\begin{aligned}\widetilde B(u,v):=\int_\Omega&\sum_{i=1}^n a_i(\bm{x},u(\bm{x}),\nabla u(\bm{x}))\dfrac{\partial v(\bm{x})}{\partial x_i}\\&+a_0(\bm{x},u(\bm{x}),\nabla u(\bm{x}))v(\bm{x})\,\mathrm{d}\bm{x}.\end{aligned}
\end{equation}
We make the following hypotheses on the functions
\begin{equation}\label{eq:cv-pde-coefficient-functions-domain}
\Omega\times\mathbb{R}\times\mathbb{R}^n\ni(\bm{x},\eta,\bm{x}ii)\mapsto a_i(\bm{x},\eta,\bm{x}ii)\in\mathbb{R},\quad i=0,\ldots,n.
\end{equation}


\begin{enumerate}[label=HP-\arabic*.,ref=HP-\arabic*]\item\label{hyp:cv-pde-spatial-continuity} For each $i=0,\ldots,n$,
\begin{equation}\label{eq:cv-pde-coefficient-spatial-continuity}
\bm{x}\mapsto a_i(\bm{x},\eta,\bm{x}ii)
\end{equation}
is in $C_b(\overline\Omega)$ for every fixed $(\eta,\bm{x}ii)\in\mathbb{R}^{n+1}$.

\item\label{hyp:cv-pde-value-continuity} For each $i=0,\ldots,n$,
\begin{equation}\label{eq:cv-pde-coefficient-value-gradient-continuity}
(\eta,\bm{x}ii)\mapsto a_i(\bm{x},\eta,\bm{x}ii)
\end{equation}
is in $C(\mathbb{R}^{n+1})$ for every $\bm{x}\in\Omega$.

\item\label{hyp:cv-pde-growth} There exists $p\in(1,\infty)$, a constant $c_0>0$, a function $k\in L^q(\Omega)$ ($\frac{1}{p}+\frac{1}{q}=1$) such that for every $\bm{x}\in\Omega$ and every $(\eta,\bm{x}ii)\in\mathbb{R}^{n+1}$ we have
\begin{equation}\label{eq:cv-pde-coefficient-growth-bound}
|a_i(\bm{x},\eta,\bm{x}ii)|\leq c[|\eta|^{p-1}+|\bm{x}ii|^{p-1}+k(\bm{x})],
\end{equation}
for each $i=0,\ldots,n$.

\item\label{hyp:cv-pde-strict-gradient-monotonicity} For every $\bm{x}ii\in\mathbb{R}^n$ and $\bm{x}ii^*\in\mathbb{R}^n$ such that $\bm{x}ii\neq\bm{x}ii^*$, and every $\eta\in\mathbb{R}$ and $\bm{x}\in\Omega$ we have
\begin{equation}\label{eq:cv-pde-strict-highest-order-monotonicity}
\sum_{i=1}^n[a_i(\bm{x},\eta,\bm{x}ii)-a_i(\bm{x},\eta,\bm{x}ii^*)](\bm{x}i_i-\bm{x}i_i^*)>0.
\end{equation}


\item\label{hyp:cv-pde-form-coercivity}
\begin{equation}\label{eq:cv-pde-absolute-bivariate-coercivity-hypothesis}
\dfrac{|\widetilde B(v,v)|}{\|v\|_{1,p}}\to\infty\quad\text{as }\|v\|_{1,p}\to\infty.
\end{equation}


\item\label{hyp:cv-pde-uniform-gradient-coercivity} For every $\bm{x}\in\Omega$ and uniformly for $|\eta|$ in bounded sets, we have
\begin{equation}\label{eq:cv-pde-uniform-gradient-coercivity-hypothesis}
\sum_{i=1}^na_i(\bm{x},\eta,\bm{x}ii)\bm{x}i_i\dfrac{1}{|\bm{x}ii|+|\bm{x}ii|^{p-1}}\to\infty\quad\text{as }|\bm{x}ii|\to\infty.
\end{equation}
\end{enumerate}

Note that by hypotheses \ref{hyp:cv-pde-spatial-continuity}, \ref{hyp:cv-pde-value-continuity} and \ref{hyp:cv-pde-growth}, the bivariate form $\widetilde B(u,v)$ is well defined for $u,v$ in $W^{1,p}(\Omega)$, where $p\in(1,\infty)$ is given as in \ref{hyp:cv-pde-growth}. Let $f\in W^{-1,q}(\Omega)$. As above, we say that $u\in W_0^{1,p}(\Omega)$ is a weak solution of the Dirichlet problem for \eqref{eq:cv-quasilinear-dirichlet-equation} if
\begin{equation}\label{eq:cv-quasilinear-dirichlet-weak-solution}
\widetilde B(u,v)=(f,v)\quad\text{for every }v\in W_0^{1,p}(\Omega).
\end{equation}
Hypotheses \ref{hyp:cv-pde-spatial-continuity}, \ref{hyp:cv-pde-value-continuity} and \ref{hyp:cv-pde-growth} also imply that for each fixed $u\in W_0^{1,p}(\Omega)$, the mapping
\begin{equation}\label{eq:cv-quasilinear-bounded-test-functional}
W_0^{1,p}(\Omega)\ni w\mapsto\widetilde B(u,w)\in\mathbb{R}
\end{equation}
is a bounded linear functional. Thus, there exists a mapping
\begin{equation}\label{eq:cv-quasilinear-induced-dual-mapping}
W_0^{1,p}(\Omega)\ni u\mapsto\overline T(u)\in W^{-1,q}(\Omega)
\end{equation}
such that
\begin{equation}\label{eq:cv-quasilinear-induced-pairing}
\widetilde B(u,w)=(\overline T(u),w)\quad\text{for all }w\in W_0^{1,p}(\Omega).
\end{equation}


We now show the following.

\begin{Theorem}
\label{thm:cv-pde-operator-is-variations-type}
Let $a_i$ satisfy hypotheses \ref{hyp:cv-pde-spatial-continuity}--\ref{hyp:cv-pde-uniform-gradient-coercivity}. Then the operator $\overline T$ is of the calculus of variations type.
\end{Theorem}

\begingroup
\renewcommand{\endProof}{\par}
\begin{Proof}
Of course, one of our goals here is to separate the effects of higher and lower derivatives; thus we define
\begin{equation}\label{eq:cv-quasilinear-three-argument-form-splitting}
\widetilde B(u,v,w):=\mathcal{B}_1(u,v,w)+\mathcal{B}_0(u,w),
\end{equation}
where
\begin{equation}\label{eq:cv-quasilinear-highest-order-three-argument-form}
\mathcal{B}_1(u,v,w):=\sum_{i=1}^n\int_\Omega a_i(\bm{x},u(\bm{x}),\nabla v(\bm{x}))\dfrac{\partial w}{\partial x_i}(\bm{x})\,\mathrm{d}\bm{x},
\end{equation}

\begin{equation}\label{eq:cv-quasilinear-lower-order-form}
\mathcal{B}_0(u,w):=\int_\Omega a_0(\bm{x},u(\bm{x}),\nabla u(\bm{x}))w(\bm{x})\,\mathrm{d}\bm{x}.
\end{equation}
Using the same argument as above, we see that there exists a mapping
\begin{equation}\label{eq:cv-quasilinear-two-argument-dual-mapping}
W_0^{1,p}(\Omega)\times W_0^{1,p}(\Omega)\ni(u,v)\mapsto\widehat T(u,v)\in W^{-1,q}(\Omega)
\end{equation}
such that
\begin{equation}\label{eq:cv-quasilinear-three-argument-duality-pairing}
\widetilde B(u,v,w)=(\widehat T(u,v),w)\quad\text{for all }w\in W_0^{1,p}(\Omega).
\end{equation}
Furthermore, we have
\begin{equation}\label{eq:cv-quasilinear-diagonal-operator-representation}
\overline T(u)=\widehat T(u,u).
\end{equation}
We now note that our results on Nemytskii operators and hypotheses \ref{hyp:cv-pde-spatial-continuity}, \ref{hyp:cv-pde-value-continuity} and \ref{hyp:cv-pde-growth} immediately imply the following lemma.
\end{Proof}
\endgroup

\begin{Lemma}
\label{lem:cv-pde-nemytskii-partial-argument-continuity}
For each $i=0,\ldots,n$ the mapping
\begin{equation}\label{eq:cv-pde-nemytskii-full-argument-continuity}
W^{1,p}(\Omega)\ni u\mapsto a_i(\bm{x},u(\bm{x}),\nabla u(\bm{x}))\in L^q(\Omega)
\end{equation}
is bounded and continuous. Furthermore, for fixed $v\in L^p(\Omega)$ the mapping
\begin{equation}\label{eq:cv-pde-nemytskii-fixed-value-continuity}
W^{1,p}(\Omega)\ni u\mapsto a_i(\bm{x},v(\bm{x}),\nabla u(\bm{x}))\in L^q(\Omega)
\end{equation}
is bounded and continuous; and for fixed $w\in W^{1,p}(\Omega)$ the mapping
\begin{equation}\label{eq:cv-pde-nemytskii-fixed-gradient-continuity}
L^p(\Omega)\ni v\mapsto a_i(\bm{x},v(\bm{x}),\nabla w(\bm{x}))\in L^q(\Omega)
\end{equation}
is bounded and continuous.
\end{Lemma}

This gives us the following corollary.

\begin{Corollary}
\label{cor:cv-pde-partial-argument-continuity}
The following hold.\begin{enumerate}\item The operator $\overline T$ is bounded and continuous.\item For each $u\in W_0^{1,p}(\Omega)$ the mapping
\begin{equation}\label{eq:cv-pde-two-argument-second-variable-continuity}
W_0^{1,p}(\Omega)\ni v\mapsto\widehat T(u,v)\in W^{-1,q}(\Omega)
\end{equation}
is bounded and continuous.\item For each $v\in W_0^{1,p}(\Omega)$ the mapping
\begin{equation}\label{eq:cv-pde-two-argument-first-variable-continuity}
W_0^{1,p}(\Omega)\ni u\mapsto\widehat T(u,v)\in W^{-1,q}(\Omega)
\end{equation}
is bounded and continuous.\end{enumerate}
\end{Corollary}

\begingroup
\renewcommand{\endProof}{\par\nobreak\hfill$\blacksquare$\par\vspace{0.1cm}}
\begingroup
\renewenvironment{Proof}{\par}{\par\nobreak\hfill$\blacksquare$\par\vspace{0.1cm}}
\begin{Proof}
This and hypothesis \ref{hyp:cv-pde-strict-gradient-monotonicity} show that conditions \ref{hyp:cv-diagonal-monotonicity} and \ref{hyp:cv-first-variable-continuity} are satisfied. Thus, to show that $\overline T$ is of the calculus of variations type and complete the proof of Theorem~\ref{thm:cv-pde-operator-is-variations-type}, we need only verify that conditions \ref{hyp:cv-vanishing-defect-weak-continuity} and \ref{hyp:cv-weak-duality-pairing} are satisfied. To check condition \ref{hyp:cv-vanishing-defect-weak-continuity} we assume that $u_j\rightharpoonup\overline u$ in $W_0^{1,p}(\Omega)$ and that
\begin{equation}\label{eq:cv-pde-vanishing-gradient-monotonicity-defect}
\begin{aligned}
&(\widehat T(u_j,u_j)-\widehat T(u_j,\overline u),u_j-\overline u)\\
&\quad=\mathcal{B}_1(u_j,u_j,u_j-\overline u)-\mathcal{B}_1(u_j,\overline u,u_j-\overline u)\\
&\quad=\int_\Omega\sum_{i=1}^n[a_i(\bm{x},u_j,\nabla u_j)-a_i(\bm{x},u_j,\nabla\overline u)]\dfrac{\partial}{\partial x_i}(u_j-\overline u)\,\mathrm{d}\bm{x}\\
&\quad\to0,
\end{aligned}
\end{equation}
and we need to show that
\begin{equation}\label{eq:cv-pde-source-unhatted-test-operator-limit}
\begin{aligned}(T(u_j,v),w)&=\mathcal{B}_1(u_j,v,w)+\mathcal{B}_0(u_j,w)\\&\to\mathcal{B}_1(\overline u,v,w)+\mathcal{B}_0(\overline u,w)\\&=(T(\overline u,v),w)\end{aligned}
\end{equation}
for all $w\in W_0^{1,p}(\Omega)$. Now by compact imbedding $u_j\to u$ (strongly) in $L^p(\Omega)$. Thus, even without using hypothesis \eqref{eq:cv-pde-vanishing-gradient-monotonicity-defect}, we have\[\begin{aligned}
\mathcal{B}_1(u_j,v,w)&=\int_\Omega\sum_{i=1}^na_i(\bm{x},u_j(\bm{x}),\nabla v(\bm{x}))\dfrac{\partial w}{\partial x_i}(\bm{x})\,\mathrm{d}\bm{x}\\
&\to\mathcal{B}_1(\overline u,v,w).
\end{aligned}\]We use the results on Nemytskii operators here. Thus, we need only show that
\begin{equation}\label{eq:cv-pde-lower-order-limit-to-verify}
\int_\Omega a_0(\bm{x},u_j(\bm{x}),\nabla u_j(\bm{x}))w(\bm{x})\,\mathrm{d}\bm{x}\to\int_\Omega a_0(\bm{x},\overline u(\bm{x}),\nabla\overline u(\bm{x}))w(\bm{x})\,\mathrm{d}\bm{x}.
\end{equation}
The proof of this in the general case is an exercise in measure theory which we shall skip. We refer the reader to \cite[Lemma 2.2, p.~184]{Li}. In our example below we use a lower-order term of the form
\begin{equation}\label{eq:cv-pde-linear-drift-p-power-lower-order-term}
a_0(\bm{x},u(\bm{x}),\nabla u(\bm{x}))=\bm{b}_1(\bm{x})\cdot\nabla u(\bm{x})+\alpha_0(\bm{x})|u(\bm{x})|^{p-2}u(\bm{x}),
\end{equation}
where $\bm{b}_1$ and $\alpha_0$ are bounded continuous functions. For such a term, condition \eqref{eq:cv-pde-lower-order-limit-to-verify} can be verified directly (that is, we need not use \eqref{eq:cv-pde-vanishing-gradient-monotonicity-defect}, though this condition is essential in the general proof).

To check condition \ref{hyp:cv-weak-duality-pairing} we assume that $u_j\rightharpoonup\overline u$ in $W_0^{1,p}(\Omega)$ and $\widehat T(u_j,v)\rightharpoonup\psi$ in $W^{-1,q}(\Omega)$; i.e.,
\begin{equation}\label{eq:cv-pde-weak-test-image-limit-assumption}
(\widehat T(u_j,v),w)=\mathcal{B}_1(u_j,v,w)+\mathcal{B}_0(u_j,w)\to(\psi,w)
\end{equation}
for every $w\in W_0^{1,p}(\Omega)$. We need to show that
\begin{equation}\label{eq:cv-pde-weak-duality-pairing-limit-goal}
(\widehat T(u_j,v),u_j)=\mathcal{B}_1(u_j,v,u_j)+\mathcal{B}_0(u_j,u_j)\to(\psi,\overline u).
\end{equation}
To do this, we write
\begin{equation}\label{eq:cv-pde-weak-duality-pairing-splitting}
(\widehat T(u_j,v),u_j)=(\widehat T(u_j,v),\overline u)+(\widehat T(u_j,v),u_j-\overline u).
\end{equation}
Thus, by \eqref{eq:cv-pde-weak-test-image-limit-assumption} we need only show that
\begin{equation}\label{eq:cv-pde-vanishing-weak-test-pairing-exercise}
\lim_{j\to\infty}(\widehat T(u_j,v),u_j-\overline u)=0.
\end{equation}
The proof of this is left to the reader (Problem~\ref{ex:cv-pde-vanishing-weak-test-pairing}). This shows that our operator is of calculus of variations type and completes the proof.
\end{Proof}
\endgroup
\endgroup

There is one more lemma to prove to get our main result on partial differential equations.

\begin{Lemma}
\label{lem:cv-pde-induced-operator-coercivity}
The operator $\overline T$ is coercive.
\end{Lemma}

The proof is left to the reader (Problem~\ref{ex:cv-pde-induced-operator-coercivity-proof}). The culmination of the previous results is the following existence theorem for quasilinear elliptic partial differential equations.

\begin{Theorem}
\label{thm:cv-quasilinear-dirichlet-existence}
Let $a_i$ satisfy hypotheses (\ref{hyp:cv-pde-spatial-continuity}--\ref{hyp:cv-pde-uniform-gradient-coercivity}), and let $p$ and $q$ be as defined in hypothesis \ref{hyp:cv-pde-growth}. Then for every $f\in W^{-1,q}(\Omega)$, there exists a weak solution $u\in W_0^{1,p}(\Omega)$ of the Dirichlet problem for the quasilinear PDE \eqref{eq:cv-quasilinear-dirichlet-equation}.
\end{Theorem}

\begin{Example}
\label{ex:cv-quasilinear-anisotropic-drift-operator}
Consider the second-order nonlinear partial differential operator
\begin{equation}\label{eq:cv-quasilinear-anisotropic-drift-example}
\begin{aligned}\widetilde A(u):=&-\sum_{i=1}^n\dfrac{\partial}{\partial x_i}\left(\left|\dfrac{\partial u}{\partial x_i}\right|^{p-2}\dfrac{\partial u}{\partial x_i}\right)\\&+\bm{b}_1(\bm{x})\cdot\nabla u(\bm{x})+\alpha_0(\bm{x})|u(\bm{x})|^{p-2}u(\bm{x}),\end{aligned}
\end{equation}
where $p\in(1,\infty)$, and where $\bm{b}_1$ and $\alpha_0$ are bounded and continuous. Note that this is the same as the operator defined in \eqref{eq:anisotropic-p-power-divergence-example} except for the difference in the lower-order terms. The reader is asked to show that there exists a constant $\widetilde C>0$ such that if
\begin{equation}\label{eq:cv-quasilinear-negative-zero-order-coefficient-bound}
\alpha_0(\bm{x})>-\widetilde C,\quad\bm{x}\in\Omega,
\end{equation}
then hypotheses \ref{hyp:cv-pde-spatial-continuity}--\ref{hyp:cv-pde-uniform-gradient-coercivity} hold (Problem~\ref{ex:cv-pde-drift-example-hypotheses}). By Theorem~\ref{thm:cv-quasilinear-dirichlet-existence} we have the following existence result.
\end{Example}

\begin{Theorem}
\label{thm:cv-quasilinear-anisotropic-drift-existence}
Let the nonlinear second-order partial differential operator $\widetilde A$ be defined by \eqref{eq:cv-quasilinear-anisotropic-drift-example}. Then for every $f\in W^{-1,q}(\Omega)$ there exists a weak solution $u\in W_0^{1,p}$ of the equation
\begin{equation}\label{eq:cv-quasilinear-anisotropic-drift-existence-equation}
\widetilde A(u)=f.
\end{equation}

\end{Theorem}

\Needspace{5\baselineskip}\textbf{Problems}

\begin{enumerate}[label=\thesection.\arabic*.,ref=\thesection.\arabic*,start=16,series=book-problems]

\item\label{ex:hemicontinuity-not-continuity-origin} We say that a mapping $T:X\to X^*$ is \textbf{hemicontinuous} at $u\in X$ if
\begin{equation}\label{eq:hemicontinuity-scalar-pairing-along-line}
\mathbb{R}\ni t\mapsto(T(u+tv),w)\in\mathbb{R}
\end{equation}
is continuous for every $v,w\in X$.\par Find a function $f:\mathbb{R}^2\to\mathbb{R}^2$ which is hemicontinuous at the origin but not continuous.

\item\label{ex:pseudo-monotone-surjectivity-proof} Prove Theorem~\ref{thm:pseudo-monotone-surjectivity}.

\item\label{ex:browder-minty-hemicontinuity-surjectivity} Show that Theorem~\ref{thm:browder-minty-surjectivity} still holds if the hypothesis of continuity is replaced by hemicontinuity.

\item\label{ex:pseudo-monotone-hemicontinuity-surjectivity} Show that Theorem~\ref{thm:pseudo-monotone-surjectivity} still holds if the hypothesis of continuity is replaced by hemicontinuity.

\item\label{ex:bounded-monotone-implies-pseudo-monotone} Show that a bounded, monotone operator is pseudo-monotone.

\item\label{ex:pseudo-monotone-strong-to-weak-continuity} Show that if $T$ is a pseudo-monotone operator and $u_j\to\overline u$ (strongly) in $X$, then $T(u_j)\rightharpoonup T(\overline u)$ in $X^*$.

\item\label{ex:cv-type-operator-hemicontinuity} Show that an operator of the calculus of variations type is hemicontinuous. (Thus, by Problem~\ref{ex:pseudo-monotone-hemicontinuity-surjectivity}, we can drop the hypothesis of continuity in Corollary~\ref{cor:cv-type-surjectivity} and the conclusion still holds.)

\item\label{ex:cv-pde-vanishing-weak-test-pairing} Assume $u_j\rightharpoonup\overline u$ in $W_0^{1,p}(\Omega)$ and $\widehat T(u_j,v)\rightharpoonup\psi$ in $W^{-1,q}(\Omega)$. Verify \eqref{eq:cv-pde-vanishing-weak-test-pairing-exercise}.

\item\label{ex:cv-pde-induced-operator-coercivity-proof} Prove Lemma~\ref{lem:cv-pde-induced-operator-coercivity}.

\item\label{ex:cv-pde-drift-example-hypotheses} Show that there is a $\widetilde C>0$ such that if \eqref{eq:cv-quasilinear-negative-zero-order-coefficient-bound} holds, then hypotheses \ref{hyp:cv-pde-spatial-continuity}--\ref{hyp:cv-pde-uniform-gradient-coercivity} are satisfied for the quasilinear differential operator $\widetilde A$ defined in \eqref{eq:cv-quasilinear-anisotropic-drift-example}. Identify which of the hypotheses \ref{hyp:monotone-pde-measurability}--\ref{hyp:monotone-pde-growth} do \emph{not} hold for this operator.

\end{enumerate}
